% Cabecera LaTeX para la exportación a PDF (xelatex + TeX Gyre Termes).
% Mapea los superíndices Unicode a \textsuperscript porque la fuente Times
% (TeX Gyre Termes) no trae esos glifos; así el texto conserva la tipografía
% académica sin perder caracteres (p. ej. 10⁻²¹, Na⁺/K⁺).
\usepackage{newunicodechar}
\newunicodechar{⁰}{\textsuperscript{0}}
\newunicodechar{¹}{\textsuperscript{1}}
\newunicodechar{²}{\textsuperscript{2}}
\newunicodechar{³}{\textsuperscript{3}}
\newunicodechar{⁴}{\textsuperscript{4}}
\newunicodechar{⁵}{\textsuperscript{5}}
\newunicodechar{⁶}{\textsuperscript{6}}
\newunicodechar{⁷}{\textsuperscript{7}}
\newunicodechar{⁸}{\textsuperscript{8}}
\newunicodechar{⁹}{\textsuperscript{9}}
\newunicodechar{⁺}{\textsuperscript{+}}
\newunicodechar{⁻}{\textsuperscript{-}}
\newunicodechar{ⁿ}{\textsuperscript{n}}
\newunicodechar{↔}{\ensuremath{\leftrightarrow}}
% Tipografía fina (mejor justificación y kerning): sello de documento académico.
\usepackage{microtype}
